% ============================================================================
% preamble.tex -- document class, thesis metadata and packages
%
% Shared by the two documents that are built from these sources:
%   thesis.tex   -- the complete thesis
%   chapter.tex  -- a single chapter, for review copies (see `make chapters')
% Both \input this file as their very first line, so the class options (in
% particular the draft switch) are set in exactly one place.
%
% The typesetting comes from the institute template at
% ~/latex-template-thesis (Fraunhofer IOSB / KIT IES,
% gitlab.cc-asp.fraunhofer.de:fuzzing/student-work/theses/latex-template-thesis).
% This file plays the role of the template's main.tex preamble block plus its
% preamble.tex; preamble/01..05 are the template's own files, kept as close to
% upstream as the switch to English allows, and preamble/06-thesis.tex holds
% what this thesis needs on top. Deviations from the template are marked
% ``DEVIATION'' where they occur.
% ============================================================================

%% ---------------------------------
%% | Draft or final                |
%% ---------------------------------
%% draft -- shows \todo notes and marks overfull boxes; use while writing
%% final -- hides them; switch before submission
%%
%% One switch drives both: the class option that marks overfull boxes and the
%% mode todonotes is loaded in (preamble/06-thesis.tex). KOMA's draft=true is a
%% key-value option that the ifdraft package cannot see, hence this boolean.
\newif\ifthesisdraft
\thesisdrafttrue

\ifthesisdraft
  \PassOptionsToClass{draft=true}{scrbook}
\else
  \PassOptionsToClass{draft=false}{scrbook}
\fi

\PassOptionsToPackage{table,dvipsnames,svgnames}{xcolor}

\documentclass[
	a4paper,
	fontsize=11pt,
	version=last,
	BCOR=10mm,
	DIV=11,
	toc=bibliography,
	toc=listof,
	toc=index,
	hyperref,
	amsmath,
	ngerman,
	twoside,
	headinclude=true,
	footnotes=multiple
 ]{scrbook}

%% --------------------------------
%% | Packages                     |
%% --------------------------------
\input{preamble/01-fonts}
% DEVIATION from the template: icomma is not loaded. It makes a comma that is
% not followed by a space a decimal separator, which is what German needs and
% what English does not: this thesis writes intervals and tuples such as
% $[0,1]$ or $(a,r,p)$, and icomma would set those without the thin space
% English typography puts after the comma.
\usepackage{fnpct}
\usepackage[final]{microtype}

% DEVIATION from the template: the thesis is written in English, so American
% English is the main language and German is only the secondary one, needed for
% the mandatory German abstract. The template has it the other way round.
\usepackage[ngerman,main=american]{babel}
\newcommand*{\ngerman}[1]{\foreignlanguage{ngerman}{#1}}
\newcommand*{\usenglish}[1]{\foreignlanguage{american}{#1}}

\usepackage[autostyle,german=quotes,english=american]{csquotes}

\usepackage{setspace}
\setstretch{1.2}

\usepackage[headsepline,plainheadsepline,draft=false]{scrlayer-scrpage}
\pagestyle{scrheadings}
\ohead[\pagemark]{\pagemark}
\chead[]{}
\ihead[]{\headmark}
\ifoot[]{}
\cfoot[]{}
\ofoot[]{}
\setkomafont{pageheadfoot}{\normalfont\sffamily}
\setkomafont{pagenumber}{\sffamily\bfseries}
\setkomafont{descriptionlabel}{\normalfont\slshape}
\deffootnote[2em]{2.5em}{1.5em}{\thefootnotemark\ }

%set a penality to prevent Schusterjungen und Hurenkinder
\widowpenalties=4 10000 10000 150 0 

% If you have the problem that many graphics are pushed at the end of the chapter try enabling the next line.
%\raggedbottom

\KOMAoptions{DIV=last}

\input{preamble/03-graphics}
\usepackage{booktabs}
\usepackage{multirow}
\usepackage{tabulary}
\usepackage{ltxtable}
\usepackage{rotfloat}
% DEVIATION from the template, which loads subfig: the thesis uses the
% subfigure environment of subcaption (see figures/ez-mcts-simulation.tex), and
% the two packages cannot be loaded together. subcaption is the maintained one;
% subfig has been unmaintained for years and clashes with caption. Captions are
% left at the KOMA defaults the template relies on, see preamble/06-thesis.tex.
\usepackage{subcaption}
\usepackage{ntheorem}
\usepackage{makeidx}
\usepackage[pdftex,unicode,final]{hyperref}
% automake would call makeindex through \write18, which needs -shell-escape.
% latexmkrc teaches latexmk to build the glossary and the acronym list instead.
% acronym gives the acronyms a list of their own instead of mixing them into
% the glossary of terms, which is what sections/backmatter.tex prints.
\usepackage[toc,acronym]{glossaries} % must be loaded after hyperref
% capitalise/noabbrev: \cref prints ``Chapter 2'', never ``chap. 2''.
% cleveref has no `american' option --- american is a dialect of its `english'.
\usepackage[capitalise,noabbrev,ngerman,english]{cleveref}
\usepackage{url}
\usepackage{mathtools}
\usepackage{mathdots} % Scale dots according to surrounding text size
\usepackage{marginnote} % Enhanced marginpar
\usepackage{stackrel} % Enhanced stackrel
\usepackage{xfrac}
\usepackage{upquote} % Set "proper" quotes in verbatim-environments
\usepackage[final]{listings}
% DEVIATION from the template, which uses style=alphabetic and sortlocale=de:
% numeric citations sorted name-year-title, as they were before the move to
% this template and as the systems-security literature this thesis cites uses
% them. The bibliography files are registered in preamble-final-setup.tex.
\usepackage[%
  backend=biber,%
  sortlocale=en,%
  style=numeric,%
  citestyle=numeric,%
  sorting=nyt,%
  pagetracker=page,%
  citereset=section,%
  doi=false,%
  isbn=false]{biblatex}
\input{glyphtounicode}
\pdfgentounicode=1

\usepackage{scrhack}

\input{preamble/05-math}
% ============================================================================
% preamble/06-thesis.tex -- what this thesis needs on top of the template
%
% preamble/01..05 are the institute template's own preamble files and are kept
% as close to upstream as the switch to English allows, so that they stay
% diffable against ~/latex-template-thesis. Everything this thesis adds lives
% here.
% ============================================================================

% ----- Figures --------------------------------------------------------------
% Figures are \input from the chapter that discusses them and reference images
% by bare filename.
\graphicspath{{figures/}}

% The template sets no caption fonts, so KOMA's defaults apply. subcaption
% loads caption, which would otherwise take over; tell it to keep using them.
\captionsetup{font=small, labelfont=bf}
\captionsetup[sub]{font=small, labelfont=bf}

% ----- Diagrams -------------------------------------------------------------
% preamble/03-graphics.tex already loads tikz with positioning, calc, fit and
% friends; the automata drawings need these two on top.
\usetikzlibrary{arrows.meta,automata}
\tikzset{terminal/.style={double,double distance=1.2pt}}

% ----- Result plots ---------------------------------------------------------
% The figures under figures/generated/ are written by scripts/evaluate.py out of
% the campaign databases and read their numbers from figures/data/*.csv. They
% need two pgfplots libraries the template does not load: fillbetween, for the
% quantile bands every distribution is drawn as, and colormaps for the named
% maps the collapse heatmaps use.
\usepgfplotslibrary{fillbetween}
\usepgfplotslibrary{colormaps}

% The generated figures set only what distinguishes them from one another;
% what they share lives here, so restyling the results does not mean
% regenerating them. Note the axis lines: the template's default is `center',
% which runs the axes through the origin. That is right for a function plot and
% wrong for a measurement read against a non-zero baseline, which is what every
% result plot in this thesis is.
% The label anchors need setting explicitly. `axis lines=center' installs label
% styles that park the labels at the far ends of the axes, which is where they
% belong when the axes cross at the origin; because the template sets that
% globally, the styles survive a per-axis switch to `axis lines=left' and the
% labels end up hanging off the corners of the plot. Re-centring them on the
% tick label band is what `axis lines=left' would have done on its own.
%
% The `rotate=90' on the ordinate label is load-bearing, not decoration: an
% unrotated label is as wide as its text, pgfplots does not count that width
% against the axis `width', and three panels of it overflow the text block.
% `scale only axis' (set by the template) makes `width' the width of the plot
% area alone, so labels, tick labels and colour bars stick out beyond it and a
% picture sized to its container overflows. Turning it off for these figures
% makes `width' the width of the whole picture, which is what a generated
% figure sized in \linewidth needs in order to fit its column.
\pgfplotsset{
  thesis result plot/.style={%
    axis lines=left,%
    axis on top=false,%
    scale only axis=false,%
    tick align=outside,%
    grid=major,%
    grid style={KITblack15,very thin},%
    legend cell align=left,%
    xlabel style={at={(ticklabel cs:0.5)},anchor=near ticklabel},%
    ylabel style={at={(ticklabel cs:0.5)},rotate=90,anchor=near ticklabel}},%
  thesis matrix plot/.style={%
    axis lines=left,%
    scale only axis=false,%
    tick align=outside,%
    enlargelimits=false,%
    axis on top,%
    xlabel style={at={(ticklabel cs:0.5)},anchor=near ticklabel},%
    ylabel style={at={(ticklabel cs:0.5)},rotate=90,anchor=near ticklabel}},%
}

% ----- Tables ---------------------------------------------------------------
% tabularx also arrives via the template's ltxtable; loading it by name states
% the dependency rather than relying on that.
\usepackage{tabularx}

% ----- Lists ----------------------------------------------------------------
% Provided by the previous template (sdqthesis.cls). The research questions in
% sections/introduction.tex use its label/ref keys, which plain enumerate has
% no equivalent for.
\usepackage{enumitem}

% ----- Units & numbers ------------------------------------------------------
\usepackage{siunitx}

% ----- Conditionals ---------------------------------------------------------
% \ifthenelse, used by the title page (an advisor line that is only printed
% when there is a second advisor) and by chapter.tex.
\usepackage{ifthen}

% ----- Algorithms -----------------------------------------------------------
% float, which algorithm needs for [H], is already loaded by the template's
% rotfloat.
\usepackage{algorithm}
\usepackage{algpseudocode}

% ----- Text macros ----------------------------------------------------------
% Provided by the previous template (sdqthesis.cls) and used throughout the
% text: \code for identifiers and paths, \model for model element names.
\newcommand{\code}[1]{\texttt{\hyphenchar\font45\relax #1}}
\newcommand{\model}[1]{\textsf{#1}}

% ----- Network names --------------------------------------------------------
% The EfficientZero networks follow MuZero's notation --- h (representation),
% g (dynamics), p (policy), v (value), vp (value prefix) --- with the parameter
% vector as an explicit argument, v(s;\theta), rather than a subscript. The
% multi-letter names need \mathit: plain math mode would set them as a product
% of single-letter variables, with the wrong spacing.
\newcommand{\vpnet}{\mathit{vp}}
\newcommand{\projnet}{\mathit{proj}}
\newcommand{\prednet}{\mathit{pred}}

% ----- Draft aids -----------------------------------------------------------
% Follows the draft switch set in preamble.tex: notes are shown while writing
% and hidden in the final document. The margin this template leaves is a little
% under todonotes' 2cm minimum. Widening \marginparwidth changes nothing in the
% final document --- \marginpar is only ever used by todonotes, which is
% disabled there.
\setlength{\marginparwidth}{2cm}
\ifthesisdraft
  \usepackage[colorinlistoftodos]{todonotes}
\else
  \usepackage[disable]{todonotes}
\fi


%% ---------------------------------
%% | Information about the thesis  |
%% ---------------------------------
%% These are plain macros, not class commands: the template builds its title
%% page (sections/titlepage.tex) and its declaration (sections/declaration.tex)
%% from them directly.

\newcommand{\worktitle}{Applications of Reinforcement Learning to Stateful
Emulator-Based Greybox Fuzzing}
\newcommand{\typeofthesis}{Master's Thesis}
\newcommand{\nameofauthor}{Bastian Engel}
\newcommand{\placeofexam}{Karlsruhe}
%% Date on the title page and the declaration: the submission date.
\newcommand{\dateofexam}{13. August 2026}
%% Start and end of the editing period, printed on the title page.
\newcommand{\editingstart}{13. February 2026}
\newcommand{\editingend}{13. August 2026}
%% The professor who grades the thesis, then the advisors. Leave \advisorB
%% empty if there is only one advisor; the title page then omits the line.
\newcommand{\reviewer}{Prof. Dr.-Ing. Jürgen Beyerer}
\newcommand{\advisorA}{Mark Leon Giraud}
\newcommand{\advisorB}{}

%% --------------------------------
%% | Setup that needs the metadata |
%% --------------------------------
% ============================================================================
% preamble-final-setup.tex -- everything that has to come after the metadata
%
% From the institute template (~/latex-template-thesis), translated to English
% and adjusted where this thesis differs; the deviations are marked.
% ============================================================================

\hypersetup{pdftitle={\worktitle},
  pdfsubject={\typeofthesis},
  pdfauthor={\nameofauthor},
  pdfstartview=Fit,
  pdfpagemode={UseOutlines},
  bookmarksnumbered=true,
  bookmarksopen=false,
  pdfborder={0 0 0},
  pdfpagelayout=TwoColumnRight
}

% ----- Theorem environments -------------------------------------------------
% ntheorem, loaded by the template. Statements are set in italics, definitions
% and examples upright.
\newtheorem{theorem}{Theorem}[chapter]
\newtheorem{lemma}[theorem]{Lemma}
\newtheorem{corollary}[theorem]{Corollary}
\newtheorem{proposition}[theorem]{Proposition}
{\theorembodyfont{\upshape}
  \newtheorem{definition}[theorem]{Definition}
  \newtheorem{example}[theorem]{Example}
}

% ----- Cross-reference names ------------------------------------------------
% The template names chapters, sections, figures and tables in German; in an
% English document cleveref's own defaults are what is wanted, so only the
% theorem environments are declared here. Plurals are what cleveref cannot
% guess.
\crefname{theorem}{Theorem}{Theorems}
\crefname{definition}{Definition}{Definitions}
\crefname{lemma}{Lemma}{Lemmata}
\crefname{corollary}{Corollary}{Corollaries}
\crefname{proposition}{Proposition}{Propositions}
\crefname{example}{Example}{Examples}
\crefname{algorithm}{Algorithm}{Algorithms}

% ----- Listings -------------------------------------------------------------
% DEVIATION from the template: Rust rather than Java, the language both
% companion repositories are written in. listings ships no Rust definition, so
% here is one --- keywords first, then the types and constructors that are
% worth setting apart.
\lstdefinelanguage{Rust}{
  keywords={as, async, await, break, const, continue, crate, dyn, else, enum,
    extern, false, fn, for, if, impl, in, let, loop, match, mod, move, mut,
    pub, ref, return, self, Self, static, struct, super, trait, true, type,
    union, unsafe, use, where, while},
  keywords=[2]{bool, char, f32, f64, i8, i16, i32, i64, i128, isize, str, u8,
    u16, u32, u64, u128, usize, Box, Option, Result, String, Vec, Some, None,
    Ok, Err},
  sensitive=true,
  comment=[l]{//},
  morecomment=[s]{/*}{*/},
  morestring=[b]{"},
}

\lstloadlanguages{
  C,
  [ISO]C++,
  [LaTeX]TeX
}

\lstset{
  basicstyle=\small\ttfamily,
  keywordstyle=\bfseries,
  commentstyle=\normalfont,
  numbers=left,
  numberstyle=\tiny,
  stepnumber=1,
  numbersep=5pt,
  tabsize=2,
  extendedchars=true,
  breaklines=true,
  showspaces=false,
  showtabs=false,
  showstringspaces=false,
  captionpos=b
}
\lstdefinestyle{nonumbers}{numbers=none}

\lstnewenvironment{latex}[1][]{%
  \lstset{%
    language=[LaTeX]TeX,%
    morekeywords={cref,%
      includegraphics,%
      toprule,%
      midrule,%
      bottomrule,%
      tikzsetnextfilename,%
      enquote,%
      newacronym,%
      longnewglossaryentry,%
      subfloat},%
    #1}%
}{}

\lstnewenvironment{rust}[1][]{%
  \lstset{language=Rust,#1}%
}{}

\lstnewenvironment{ccode}[1][]{%
  \lstset{language=C,#1}%
}{}

% ----- Margin notes ---------------------------------------------------------
\renewcommand*{\marginfont}{\raggedright\sffamily\footnotesize}
\let\omarginnote\marginnote
\renewcommand{\marginnote}[1]{\leavevmode\omarginnote{#1}\ignorespaces}

% ----- Index, glossary, acronyms --------------------------------------------
% Built by latexmk, see latexmkrc.
\makeindex

\setacronymstyle{long-short}
\makeglossaries

% ----- Bibliography ---------------------------------------------------------
% Single bibliography, kept in sync with Zotero by the editor: citing a key that
% is not in here yet appends its entry automatically (:ZoteroSync adds any that
% were pasted in by hand).
\addbibresource{references/library.bib}
% \nocite{*} % activate to make latex print all literature, even the ones that were not cited in the thesis


%% Acronyms and glossary entries, printed by sections/backmatter.tex.
%!TEX root = ../thesis.tex
% Acronym definitions, printed as the list of acronyms by
% sections/backmatter.tex. Inputted from preamble.tex, since \newacronym has to
% run in the preamble.
%
% Use them in the text with \gls{mcts} (long at first use, short afterwards),
% \glspl{mcts} for the plural, \acrlong{mcts} / \acrshort{mcts} to force one
% form. An acronym that is never used in the text is still printed in the list:
% sections/backmatter.tex calls \glsaddallunused.
%
% TODO: MOGI is the name of the emulation framework, not an acronym with a
% documented expansion --- add it here once the MOGI paper says what it stands
% for.

% ----- Fuzzing --------------------------------------------------------------
\newacronym{cgf}{CGF}{coverage-guided fuzzing}
\newacronym[longplural={systems under test}, shortplural={SUTs}]
  {sut}{SUT}{system under test}
\newacronym{afl}{AFL}{American Fuzzy Lop}
\newacronym{bb}{BB}{basic block}

% ----- Reinforcement learning -----------------------------------------------
\newacronym{rl}{RL}{reinforcement learning}
\newacronym[longplural={Markov decision processes}, shortplural={MDPs}]
  {mdp}{MDP}{Markov decision process}
\newacronym[longplural={partially observable Markov decision processes},
  shortplural={POMDPs}]
  {pomdp}{POMDP}{partially observable Markov decision process}
\newacronym{td}{TD}{temporal difference}
\newacronym{mcts}{MCTS}{Monte-Carlo tree search}
\newacronym{ucb}{UCB}{upper confidence bound}
\newacronym{uct}{UCT}{upper confidence bounds applied to trees}
\newacronym{dqn}{DQN}{deep Q-network}
\newacronym{ppo}{PPO}{proximal policy optimization}
\newacronym{sarsa}{SARSA}{state-action-reward-state-action}
\newacronym{gpi}{GPI}{generalized policy iteration}
\newacronym{per}{PER}{prioritized experience replay}
\newacronym{mlp}{MLP}{multilayer perceptron}
\newacronym{lstm}{LSTM}{long short-term memory}
\newacronym{kl}{KL}{Kullback-Leibler}

% ----- Mobile network protocols ---------------------------------------------
% Only used by the discussion of LEFT in the related-work chapter.
\newacronym{lte}{LTE}{Long Term Evolution}
\newacronym[longplural={user equipments}, shortplural={UEs}]
  {ue}{UE}{user equipment}
\newacronym{rrc}{RRC}{Radio Resource Control}
\newacronym{emm}{EMM}{EPS Mobility Management}

% ----- Systems --------------------------------------------------------------
\newacronym{cpu}{CPU}{central processing unit}
\newacronym{gpu}{GPU}{graphics processing unit}
\newacronym{mmu}{MMU}{memory management unit}
\newacronym{rng}{RNG}{random number generator}
\newacronym{prg}{PRG}{pseudo-random generator}
\newacronym{iot}{IoT}{Internet of Things}

% ----- Input formats and protocols ------------------------------------------
\newacronym{json}{JSON}{JavaScript Object Notation}
\newacronym{xml}{XML}{Extensible Markup Language}
\newacronym{http}{HTTP}{Hypertext Transfer Protocol}
\newacronym{opcua}{OPC~UA}{Open Platform Communications Unified Architecture}

